\originalchapter{8.513作业5--8：量子临界与一维异常}
\psource{Xiao-Gang Wen，8.513 Fall 2021，公开作业5--8。}
\originalsection{作业5}
\begin{exercise}[原题1(a)--(b)]
自旋1链为
\begin{equation}
H=\sum_i[-J(S_i^xS_{i+1}^x+S_i^yS_{i+1}^y)-BS_i^z+V(S_i^z)^2],\quad J>0.
\end{equation}
试验态为各站$(\phi_i,1,\widetilde\phi_i)^T/\sqrt{1+|\phi_i|^2+|\widetilde\phi_i|^2}$的张量积。(a)大$V$的Mott态$\phi=\widetilde\phi=0$附近求二次相空间作用量。(b)从大$V$降低$V$，求$B=0$和$B\ne0$首次超流失稳的低能色散与动态指数。
\end{exercise}
\begin{solution}
$\avg{S^z}=|\phi|^2-|\widetilde\phi|^2$，$\avg{(S^z)^2}=|\phi|^2+|\widetilde\phi|^2$到二阶，$\avg{S^+}=\sqrt2(\phi^*+\widetilde\phi)$到一阶。Berry项为$\ii(\phi^*\dot\phi+\widetilde\phi^*\dot{\widetilde\phi})$差总导数。二次能$\sum[(V-B)|\phi|^2+(V+B)|\widetilde\phi|^2]-J\sum[(\phi_i^*+\widetilde\phi_i)(\phi_{i+1}+\widetilde\phi_{i+1}^*)+\mathrm{c.c.}]$。

Fourier后对$(\phi_k,\widetilde\phi_{-k}^*)$的动态矩阵为$\left(\begin{smallmatrix}V-B-2J\cos k&-2J\cos k\\2J\cos k&-V-B+2J\cos k\end{smallmatrix}\right)$，正激发两支$E_\pm=\sqrt{V(V-4J\cos k)}\pm B$。$B=0$在$V_c=4J$闭隙，$E\simeq2\sqrt2J|k|$，$z=1$。$B\ne0$首次失稳满足$V_c(V_c-4J)=B^2$，$V_c=2J+\sqrt{4J^2+B^2}$，低支$E\simeq V_cJk^2/|B|$，$z=2$。这是同一变分模型的粒子--空穴对称尖端与普通边界的区别；不将一维平均场相图当精确临界参数。
\end{solution}
\begin{exercise}[原题2(a)--(b)]
一维右行复费米子$E_k=\hbar vk,k=2\pi n/L$，基态填满$k\le0$。(a)保持总费米子数的低能多体谱，给前五个激发能级简并，与右行玻色子正整数模比较。(b)求热容并比较。
\end{exercise}
\begin{solution}
把费米海中各粒子从最高占据开始依次抬升，Pauli限制使抬升量成为非增整数序列，即整数分拆。激发能$n(2\pi\hbar v/L)$，$n=0,1,2,3,4,5$简并$1,1,2,3,5,7$；若“前五级”含基态则取前五项。玻色各模占据生成函数$\prod_{m\ge1}(1-q^m)^{-1}$给同谱，零模总电荷扇区已固定。

复费米子相对海的正能粒子与负能空穴均热激发，$E/L=2\int_0^\infty(\hbar vk)f(\hbar vk)\dd k/(2\pi)=\pi(\kb T)^2/(12\hbar v)$，用$\int x/(\e^x+1)=\pi^2/12$。因此$C/L=\pi\kb^2T/(6\hbar v)$，同一支手性玻色子的结果。Majorana不是复费米子，热容另减半。
\end{solution}
\originalsection{作业6}
\begin{exercise}[原题1(a)--(f)]
开链XY模型$H=-\sum_{i=1}^{N-1}(\sigma_i^x\sigma_{i+1}^x+\sigma_i^y\sigma_{i+1}^y+h\sigma_i^z)$，旋转$\prod_i\e^{\ii\theta\sigma_i^z/2}$为$U(1)$。(a)作Jordan--Wigner映射。(b)按原题改为偶数$N$的平移不变费米环，求$E_{\rm ring}(h,N)$。(c)找$\pm h_c$。(d)小$|h|$低能线性色散、中心荷与基态$1/N$项。(e)接近$h_c$速度如何消失，临界色散是什么？(f)中子给动量$k$并使自旋改变1，态$\sum_j\e^{\ii kj}\sigma_j^-\ket0$是否像$\sum\e^{\ii kj}c_j\ket0$一样是能量态？描述散射谱。
\end{exercise}
\begin{solution}
(a)$\sigma_j^z=1-2n_j$，XY键变$-2(c_j^\dagger c_{j+1}+\mathrm{h.c.})$，磁场变$2hn_j-h$。题面将场也求和到$N-1$，最后站无场，是边界缺项；(b)的平移环按题目指定改为所有站同场，不能声称与该开链每个有限尺寸能量相同。

(b)环取所选边界动量集$\mathcal K$，$\epsilon_k=-4\cos k+2h$，精确$E_{
m ring}=-hN+\sum_{k\in\mathcal K,\epsilon_k<0}\epsilon_k$，零能态须说明占据与宇称选择。体极限$|h|<2$时$k_F=\arccos(h/2)$，$E/N=-h+[2hk_F-4\sin k_F]/\pi$；$h>2$为空费米海$E/N=-h$，$h<-2$全满$E/N=h$，所以(c)$h_c=2$。

(d)两个费米点附近$E_{k_F+q}-E_F\simeq\pm vq$，$v=4\sin k_F$。左右复费米子构成非手性$c_R=c_L=1$的理论，通常称中心荷$c=1$；$c_{\rm tot}=2$。选CFT真空扇区时$E_0=Ne_\infty-\pi v/(6N)+\cdots$。一般有限环因粒子数舍入与边界扭曲另有零模$O(1/N)$项，必须用(b)的精确和核对，不能对任意周期动量海不加扇区条件套同一系数。

(e)$v=4\sqrt{1-h^2/4}\to0$，在$h=2$带底$\epsilon_k=2k^2+O(k^4)$，$z=2$，不是保持有限光速的Lorentz理论。(f)$\sigma_j^-=(\prod_{l<j}(1-2n_l))c_j^\dagger$或其共轭，依所选升降约定。该字符串作用于半填充海会叠加许多不同粒子--空穴激发，不是单个$k$费米本征态。动态谱为$\sum_n|\avg{n|S^-_k|0}|^2\delta(\omega-E_n+E_0)$，热力学极限形成连续谱与幂律边缘；小动量低能边缘与$v|q|$关联，但不等于一条裸$\epsilon_k$的$\delta$曲线。有限链当然仍由离散峰组成，连续谱是尺寸极限。
\end{solution}
\begin{exercise}[原题2]
环上两个手性复费米场$H=\int\dd x[c_1^\dagger v_1\ii\partial_xc_1+c_2^\dagger v_2\ii\partial_xc_2+M c_1^\dagger c_2+M^*c_2^\dagger c_1]$，标准反对易。分别$v_1v_2>0$与$<0$求谱，说明质量混合何时仍无隙，联系净手性中心荷。
\end{exercise}
\begin{solution}
Fourier$\e^{\ii kx}$下矩阵$\left(\begin{smallmatrix}-v_1k&M\\M^*&-v_2k\end{smallmatrix}\right)$，根$E_\pm=-(v_1+v_2)k/2\pm\sqrt{(v_1-v_2)^2k^2/4+|M|^2}$。若$v_1v_2>0$，行列式$v_1v_2k^2-|M|^2$仍在$k=\pm|M|/\sqrt{v_1v_2}$为0，只移开交叉，无全隙；净手性$c_-=\pm2$。若$v_1v_2<0$，无实零点，最小正能$2|M|\sqrt{-v_1v_2}/|v_1-v_2|>0$，净手性0。只看$k=0$的$\pm|M|$不能判断整个谱是否有隙；连续理论需交叉仍在其紫外有效窗口。
\end{solution}
\originalsection{作业7}
\begin{exercise}[原题1(a)]
手性玻色子的$E_n=n\,2\pi\hbar v/L$简并$p(n)\sim\e^{2\pi\sqrt{n/6}}/(4\sqrt3n)$。以此微正则态密度做鞍点求热容，并推广中心荷$c$的$D(E)\sim\exp[L\sqrt{c\pi(E/L)/(3\hbar v)}]$。区分$c$是手性总态计数还是非手性左右各一份。
\end{exercise}
\begin{solution}
主指数$\log D=\sqrt{\pi LE/(3\hbar v)}$，$\log Z$积分鞍点满足$\partial_E\log D=\beta$，所以$E_*=\pi L/(12\hbar v\beta^2)$，$C=\pi L\kb^2T/(6\hbar v)$。代数前因子产生有限尺寸次阶修正，在$\kb T\gg2\pi\hbar v/L$可略。将热能系数乘总手性计数$c_{\rm tot}$再用$S=2E/T$，得到题示指数的$c=c_{\rm tot}$。如果采用非手性教科书$c_R=c_L=c$，指数中应放$2c$，不可混用。
\end{solution}
\begin{exercise}[原题2(a)--(d)]
周期横场Ising链长$L\equiv0\pmod4$，$H=-\sum(J\sigma_i^x\sigma_{i+1}^x+h\sigma_i^z)$，$J>0$。(a)以$\lambda_i^{x,y}=(\prod_{j<i}\sigma_j^z)\sigma_i^{x,y}$作映射。(b)比较周期、反周期Majorana模型与自旋$Z=\prod\sigma^z$扇区，找正确FNP投影。(c)求两种费米边界真空宇称，单独考虑$k=0,\pi$。(d)$J>|h|$求最低两自旋态劈裂的指数形式和$\xi$；题面引用附论文未随公开作业附上，但可独立分析支点。
\end{exercise}
\begin{solution}
(a)内键$\sigma_i^x\sigma_{i+1}^x=-\ii\lambda_i^y\lambda_{i+1}^x$，边界多串$Z$，所以$c_{L+1}=-Zc_1$。(b)$FNP=\prod_i(\ii\lambda_i^y\lambda_i^x)=Z$，偶自旋扇区必须取反周期费米且投偶宇称；奇自旋取周期且投奇宇称。原题“even or maybe odd”是要确定这个配对，不能把两个完整费米空间都保留。

(c)反周期$k,-k$全部成对，BCS真空偶。周期的$k=0,\pi$配对幅度零，单模能$\xi_k=2(h-J\cos k)$，负则填。$-J<h<J$仅$k=0$填，真空奇；$h>J$两者空，$h<-J$两者满，皆偶。任一所需投影不匹配时最低允许态需加最便宜的准粒子。

(d)$E(k)=2\sqrt{J^2+h^2-2Jh\cos k}$，基态能在匹配扇区为$-\frac12\sum_kE(k)$。两动量格点错半步，Poisson求和差由第$L$阶及更高Fourier系数决定。最近复支点$k=\pm\ii\kappa$（负$h$可平移实部$\pi$），$\cosh\kappa=(J^2+h^2)/(2J|h|)$，$\kappa=\log(J/|h|)$。平方根支点的Fourier系数带代数因子，故$\Delta E=O(L^\mu)\e^{-L/\xi}$，$\xi^{-1}=\log(J/|h|)$，$h=0$是严格零劈裂极限。只求指数不需取得所引论文，也没有从该论文猜造额外问题。
\end{solution}
\begin{exercise}[原题3(a)--(c)]
反周期手性Majorana环$H=\sum_{k>0}vk\lambda_k^\dagger\lambda_k$，$k=2\pi(n+1/2)/L$，$\lambda_{-k}=\lambda_k^\dagger$，$\lambda_k\ket0=0$于$k>0$。(a)求虚时间有序$\avg{T\lambda_{k_1}(\tau)\lambda_{k_2}(0)}$，注意$\tau<0$负号。(b)$\lambda(x)=L^{-1/2}\sum\lambda_k\e^{\ii kx}$，证明Hermitian与局域反对易。(c)求环相关与无穷长极限。
\end{exercise}
\begin{solution}
(a)$\tau\ge0$只有$k_1>0,k_2=-k_1<0$非零，权重$[1-\operatorname{sgn}k_2]\e^{-v|\tau k_2|}\delta_{k_1+k_2}/2$；$\tau<0$先把算符交换，权重$-[1+\operatorname{sgn}k_2]\e^{-v|\tau k_2|}\delta_{k_1+k_2}/2$。(b)共轭后换$k\to-k$得自身，反对易为$L^{-1}\sum\e^{\ii k(x-y)}$，在反周期环上是带对应边界的$\delta$核，不同点为0。

(c)$\tau>0$求正模几何和得到$G=L^{-1}/[\e^{-\ii\pi(x+\ii v\tau)/L}-\e^{\ii\pi(x+\ii v\tau)/L}]$，即第\ref{ch:chiral}章的$\ii/[2L\sin(\pi z/L)]$。$L\to\infty$得$\ii/(2\pi z)=(-2\pi\ii)^{-1}/z$。$\tau<0$以时间排序的解析延续处理，不能把正模和直接当所有$\tau$收敛的级数。
\end{solution}
\originalsection{作业8}
\begin{exercise}[原题1(a)--(c)]
$U(1)$紧致相位低能Lagrangian为$L=\int_0^L\dd x[(\partial_t\theta)^2/(4\pi V_2)-V_1(\partial_x\theta)^2/(4\pi)+\rho_0\partial_t\theta]$，$\rho_0=|\bar\phi|^2/a$，$\hbar=1$。(a)零模$\theta=2\pi mx/L+\theta_0$求每$(m,N)$扇区最低能。(b)求总玻色子数。(c)平移引起相位$2\pi m\Delta x/L$，求总动量；也可取$\theta=2\pi m[x-x_0(t)]/L+\theta_0$由共轭动量求之，解释混合异常。
\end{exercise}
\begin{solution}
零模$p_0=L\dot\theta_0/(2\pi V_2)+\rho_0L$。紧致角要求$p_0=N\in\mathbb Z$，所以$\dot\theta_0=2\pi V_2(N-\rho_0L)/L$、$E_{m,N}=\pi V_2(N-\rho_0L)^2/L+\pi V_1m^2/L$，$N$就是该扇区的粒子数。固定$m$的最低态选择最接近$\rho_0L$的整数$N$；半整数位置可简并，不能擅自当$N$连续。

平移使态获$\e^{\ii N2\pi m\Delta x/L}$，相应$P=2\pi mN/L$按该主动平移约定。用$x-x_0$的参数则$p_{x_0}=-(2\pi m/L)N$，符号来自把图样位置与场主动平移相反定义；一致恢复同物理动量。改变$U(1)$扭曲$m$泵送平移量子数，即混合异常，Berry总导数在量子零模里有实际作用。
\end{solution}
\begin{exercise}[原题2(a)--(e)]
环上$H_0=\sum_i[-t(b_i^\dagger b_{i+1}+\mathrm{h.c.})+\mu n_i+Vn_in_{i+1}]$，整数粒子数$N_0$、平移$Tb_iT^{-1}=b_{i+1}$，基态$T\ket\Phi=\e^{\ii K_0}\ket\Phi$。定义$U_\theta=\exp[(\ii\theta/L)\sum_{i=0}^{L-1}in_i]$、$\ket{\Phi_\theta}=U_\theta\ket\Phi$。(a)$\theta=2\pi r$时求动量差。(b)求$U_\theta^{-1}H_{i,i+1}U_\theta$，说明跨边界。(c)求其二阶导数阶数。(d)利用$\pm2\pi$变分能均不低于基态证明每键能差$O(L^{-2})$。(e)推总能差$O(L^{-1})$与既约填充$p/q$的$q$个近简并低能态；讨论无理填充。题页“spin1/2”与标准玻色对易同时出现不相容，证明只用整数$n_i$和$[n_i,b_j]=-\delta_{ij}b_j$，对硬核或有限密度玻色模型分别注明界。
\end{exercise}
\begin{solution}
(a)$TU_\theta T^{-1}=U_\theta\e^{-\ii\theta N_0/L}\e^{\ii\theta n_0}$，整数$n_0$且$\theta=2\pi r$时最后因子1，所以$\Delta K=-\theta N_0/L$模$2\pi$，不同动量的态正交。

(b)内键$H_i^\theta=-t(\e^{\ii\theta/L}b_i^\dagger b_{i+1}+\mathrm{h.c.})+\mu n_i+Vn_in_{i+1}$。跨边界的纯规范变换相位实际为$\e^{-\ii\theta(L-1)/L}$，仅$\theta=2\pi r$时同$\e^{\ii\theta/L}$。要取任意$\theta$导数，可把内键函数解析延拓为均匀扭曲键；不能声称任意非整数通量都是全局纯规范。

(c)局域跃迁系数二阶导数为$O(L^{-2})$，在平移基态、有限密度$\rho$下$|\avg{b_i^\dagger b_{i+1}}|\le\sqrt{\avg{n_i}\avg{n_{i+1}}}=\rho$，所以每键能$\epsilon(\theta)$的二阶导数也$O(L^{-2})$。对普通无界玻色算符不是全Hilbert空间算符范数界，题目需要的是此期望界；硬核模型可直接用局域算符有界性。

(d)在$\theta=\pm2\pi$，扭曲态仍平移本征，能$L\epsilon(\pm2\pi)$。变分原理使两差$\epsilon(\pm2\pi)-\epsilon(0)\ge0$。Taylor正负相加消去一阶，和至多$C/L^2$，故各自也至多$C/L^2$，不必假定基态电流预先为0。(e)乘$L$得$O(1/L)$总差。既约$p/q$填充的$r=0,\ldots,q-1$扭曲动量相异，固定$q$时都有消失能差，排除唯一有隙对称基态。无理填充极限可产生任意多相异低能动量；若要求有限基态简并且保持平移，必须无隙。原题的“无论相互作用都无隙”还依赖局域、有限密度、这些对称与热力学极限条件，不能扩展到非局域或无限简并的任意模型。
\end{solution}
